% -------------------------------------------------------------------------
% ANEXOS - MATERIAL COMPLEMENTARIO
% -------------------------------------------------------------------------

Este anexo recoge información complementaria que, sin ser imprescindible para la lectura del cuerpo principal, facilita la reproducibilidad del trabajo y documenta decisiones de implementación.

\section{Entorno de cómputo y reproducibilidad}

Todos los experimentos se implementaron en Python utilizando PyTorch como marco de aprendizaje profundo, junto con las bibliotecas \texttt{timm} y \texttt{transformers} para la carga de los \textit{backbones} preentrenados. El entrenamiento se realizó sobre unidades de procesamiento gráfico (GPU) con aceleración CUDA.

Para garantizar la reproducibilidad de los resultados se fijaron de forma explícita las semillas de los generadores de números aleatorios de PyTorch y NumPy, y se habilitó el modo determinista de cuDNN. Cada arquitectura se entrenó con diez semillas independientes (de la 42 a la 51), lo que permitió cuantificar la estabilidad del entrenamiento frente a la inicialización.

La selección de cortes se precalculó y almacenó en caché, de modo que el selector MobileNetV2 no debe ejecutarse en cada iteración de entrenamiento, lo que acelera notablemente el proceso y asegura que todos los modelos se entrenen sobre exactamente los mismos cortes.

\section{Configuración de aumentos de datos por plano}

Los aumentos de datos (\textit{data augmentation}) se aplicaron únicamente sobre el conjunto de entrenamiento y se ajustaron de forma independiente por plano anatómico y arquitectura, buscando un equilibrio entre regularización y preservación de la información anatómica. La Tabla~\ref{tab:augmentation} resume la intensidad de aumentos empleada.

\begin{table}[H]
\centering
\small
\caption{Intensidad de los aumentos de datos aplicados por arquitectura y plano anatómico.}
\label{tab:augmentation}
\begin{tabularx}{\textwidth}{
>{\raggedright\arraybackslash}X
>{\centering\arraybackslash}m{2.6cm}
>{\centering\arraybackslash}m{2.6cm}
>{\centering\arraybackslash}m{2.6cm}
}
\toprule
\textbf{Arquitectura} &
\textbf{Sagital} &
\textbf{Coronal} &
\textbf{Axial} \\
\midrule
ResNet50   & Agresivo      & Moderado & Agresivo \\
ViT-Small  & Conservador   & Moderado & Agresivo \\
Swin-Tiny  & Agresivo      & Agresivo & Agresivo \\
\bottomrule
\end{tabularx}
\end{table}

Los modos de aumento se definieron de forma incremental. El modo \textit{conservador} aplica únicamente rotaciones y traslaciones pequeñas que preservan la anatomía; el modo \textit{moderado} añade volteos horizontales, variaciones de brillo y contraste y transformaciones afines suaves; y el modo \textit{agresivo} incorpora además rotaciones más amplias, distorsiones de perspectiva y mayores variaciones fotométricas.

\section{Esquemas de fusión multi-vista explorados}

Además del promedio simple de probabilidades utilizado en el cuerpo principal, durante el desarrollo se exploraron esquemas alternativos de fusión multi-vista:

\begin{itemize}
    \item \textbf{Promedio simple.} Media aritmética de las probabilidades de los tres planos. Es el esquema adoptado por no introducir parámetros adicionales y permitir una comparación directa entre arquitecturas.
    \item \textbf{Votación por máximo.} Se toma la probabilidad máxima entre los tres planos, favoreciendo la detección cuando cualquiera de las vistas considera probable la lesión.
    \item \textbf{Promedio ponderado por AUC.} Cada plano contribuye con un peso proporcional a su AUC de validación, otorgando mayor influencia a las vistas más fiables.
\end{itemize}

Estas alternativas ofrecieron resultados próximos al promedio simple, sin una mejora consistente que justificara su adopción frente a la mayor simplicidad e interpretabilidad del promedio aritmético.

